\documentclass[12pt]{article}
\usepackage[utf8]{inputenc}
\usepackage[T1]{fontenc}
\usepackage[margin=1in]{geometry}
\usepackage{setspace}
\usepackage{xcolor}
\usepackage[colorlinks=true, urlcolor=blue!60!black, linkcolor=blue!60!black]{hyperref}
\setlength{\parindent}{2em}
\setlength{\parskip}{1.25em}
\renewcommand{\baselinestretch}{1.0}
\begin{document}
\begin{center}
{\LARGE \textit{Milliken v. Bradley}, 418 U.S. 717 (1974)\par}
\vspace{0.8em}
{\large Case Brief\par}
\end{center}
\vspace{1em}

\section{Memory Jogger}
In \href{https://www.courtlistener.com/opinion/109102/milliken-v-bradley/}{\textit{Milliken v. Bradley}}, the Supreme Court barred federal courts from imposing an inter-district metropolitan school desegregation remedy across suburban district lines unless plaintiffs prove an inter-district constitutional violation that produced a significant segregative effect in another district.

\section{Detailed Case Facts}
In April 1970, facing severe racial polarization in the Detroit Public Schools---which then enrolled roughly 276,000 to 290,000 students, approximately 63.6\% Black and 34.8\% white, amid a tri-county metropolitan area (Wayne, Oakland, and Macomb Counties) that was overwhelmingly white---the Detroit Board of Education voluntarily adopted the ``April 7, 1970 Plan,'' a partial high school desegregation and decentralization plan (\href{https://supreme.justia.com/cases/federal/us/418/717/}{Justia, 418 U.S. at 722--23, 729 n.10, 765 \& n.1}; \href{https://www.michbar.org/file/journal/pdf/pdf4article1911.pdf}{State Bar of Michigan}). In response, the Michigan Legislature enacted Act 48 of the Public Acts of 1970, whose Section 12 delayed and effectively nullified the April 7 Plan and mandated ``free choice'' and ``neighborhood schools'' criteria (\href{https://supreme.justia.com/cases/federal/us/418/717/}{418 U.S. at 722--23, 727 \& n.6}).

In August 1970, the Detroit Branch of the NAACP and individual Detroit parents and students (including Ronald Bradley) filed a class action on behalf of all Detroit public school children and parents against Michigan Governor William G. Milliken, the State Attorney General, the State Board of Education, the State Superintendent of Public Instruction, and the Detroit Board of Education (\href{https://supreme.justia.com/cases/federal/us/418/717/}{418 U.S. at 722}; \href{https://clearinghouse.net/case/11076/}{Civil Rights Litigation Clearinghouse}). Plaintiffs challenged Act 48 under the Fourteenth Amendment's Equal Protection Clause and sought a decree dismantling \textit{de jure} segregation in the Detroit school system (\href{https://supreme.justia.com/cases/federal/us/418/717/}{418 U.S. at 722--23}).

After a 41-day trial, the District Court found extensive \textit{de jure} segregative acts by the Detroit Board: creating optional attendance zones in racially transitioning neighborhoods that let white students escape Black schools; drawing attendance boundaries north-south rather than east-west to maximize separation; busing Black students past closer white schools with vacant seats to predominantly Black schools (while never busing white students to predominantly Black schools despite 22,961 vacant seats in schools 90\% or more Black); and siting new school construction so that 11 of 14 schools opened in 1970--1971 were over 90\% Black (\href{https://supreme.justia.com/cases/federal/us/418/717/}{418 U.S. at 724--26}). The District Court also held the State of Michigan responsible through Act 48, state oversight of school construction and site selection prior to 1962, discriminatory denial of state pupil-transportation aid to Detroit while subsidizing suburban busing, and the State's vicarious responsibility for Detroit as a state instrumentality (\href{https://supreme.justia.com/cases/federal/us/418/717/}{418 U.S. at 726--28}). Critically, however, the court took no proof that any of the 85 outlying suburban school districts had committed acts of \textit{de jure} segregation or that school district boundaries had been drawn for racial purposes (\href{https://supreme.justia.com/cases/federal/us/418/717/}{418 U.S. at 730 n.11, 733}).

\section{Procedural History}
U.S. District Judge Stephen J. Roth initially denied plaintiffs' motion for a preliminary injunction against Act 48, but the Sixth Circuit reversed, holding that Act 48 unconstitutionally thwarted the April 7 Plan, and remanded for an expedited merits trial (\href{https://supreme.justia.com/cases/federal/us/418/717/}{418 U.S. at 723}, citing 433 F.2d 897 (6th Cir. 1970)). After a second interlocutory appeal over the Board's ``Magnet Plan'' (438 F.2d 945 (6th Cir. 1971)), Judge Roth conducted a 41-day trial from April 6 to July 22, 1971, and issued his liability ruling on September 27, 1971 (338 F. Supp. 582 (E.D. Mich. 1971)), finding \textit{de jure} segregation by both the Detroit Board and State officials (\href{https://supreme.justia.com/cases/federal/us/418/717/}{418 U.S. at 723--28}).

Judge Roth deferred a motion by intervening Detroit parents to join the 85 outlying tri-county school districts as ``premature,'' yet ordered State officials to submit three-county metropolitan desegregation plans alongside Detroit-only plans (\href{https://supreme.justia.com/cases/federal/us/418/717/}{418 U.S. at 728--30}). On March 15, 1972, he granted limited intervention to outlying districts, barring them from reopening liability or prior rulings (\href{https://supreme.justia.com/cases/federal/us/418/717/}{418 U.S. at 730--31}). On March 28, 1972, Judge Roth rejected all three Detroit-only plans as incapable of accomplishing desegregation---reasoning they would leave Detroit's schools 75\% to 90\% Black and accelerate white flight---and declared that ``[s]chool district lines are simply matters of political convenience'' (\href{https://supreme.justia.com/cases/federal/us/418/717/}{418 U.S. at 732--33}). On June 14, 1972 (345 F. Supp. 914 (E.D. Mich. 1972)), he designated Detroit plus 53 suburban school districts (covering roughly 780,000 total students) as a metropolitan ``desegregation area'' divided into 15 clusters, and on July 11, 1972, ordered the State-funded purchase of at least 295 buses (\href{https://supreme.justia.com/cases/federal/us/418/717/}{418 U.S. at 733--34 \& n.14}).

Sitting en banc on June 12, 1973 (484 F.2d 215 (6th Cir. 1973)), the Sixth Circuit affirmed the liability findings and held that a metropolitan inter-district remedy was the only feasible way to desegregate Detroit and fell within the District Court's equity powers, while vacating the bus-purchase order and remanding so affected suburban districts could be joined under Federal Rule of Civil Procedure 19 (\href{https://supreme.justia.com/cases/federal/us/418/717/}{418 U.S. at 734--36}). The Supreme Court granted certiorari (414 U.S. 1038 (1973)), heard argument on February 27, 1974, and decided the consolidated cases on July 25, 1974 (\href{https://supreme.justia.com/cases/federal/us/418/717/}{418 U.S. at 717, 721 n.1}).

\section{Judicial Votes}
\textbf{Vote:} 5--4, reversing the judgment of the Sixth Circuit and remanding for prompt formulation of a Detroit-only desegregation decree (\href{https://supreme.justia.com/cases/federal/us/418/717/}{418 U.S. at 720, 753}).
\begin{itemize}
\item \textbf{Majority (5):} Chief Justice Warren E. Burger (author), joined by Justices Potter Stewart, Harry A. Blackmun, Lewis F. Powell, Jr., and William H. Rehnquist.
\item \textbf{Concurrences (1):} Justice Potter Stewart filed a concurring opinion.
\item \textbf{Dissents (4):}
\begin{itemize}
\item Justice William O. Douglas filed a dissenting opinion.
\item Justice Byron R. White filed a dissenting opinion, joined by Justices William O. Douglas, William J. Brennan, Jr., and Thurgood Marshall.
\item Justice Thurgood Marshall filed a dissenting opinion, joined by Justices William O. Douglas, William J. Brennan, Jr., and Byron R. White.
\end{itemize}
\end{itemize}

\section{Holding}
\textbf{Judgment:} Reversed and remanded to the Sixth Circuit for further proceedings leading to a prompt desegregation decree limited to the Detroit city school system (\href{https://supreme.justia.com/cases/federal/us/418/717/}{418 U.S. at 753}).

\textbf{Rule Announced:} Under the Fourteenth Amendment's Equal Protection Clause and equitable remedial principles established in \textit{Swann v. Charlotte-Mecklenburg Board of Education}, 402 U.S. 1 (1971), ``the scope of the remedy is determined by the nature and extent of the constitutional violation'' (\href{https://supreme.justia.com/cases/federal/us/418/717/}{418 U.S. at 744}). A federal court may not impose a multi-district, area-wide desegregation remedy for single-district \textit{de jure} segregation unless plaintiffs first show an \textit{inter-district constitutional violation and inter-district segregative effect}---specifically, that racially discriminatory acts of the State or of one or more local school districts were a substantial cause of inter-district segregation, or that school district boundary lines were deliberately drawn on the basis of race (\href{https://supreme.justia.com/cases/federal/us/418/717/}{418 U.S. at 744--45}). Moreover, desegregation does not require any particular racial balance in each school, grade, or classroom (\href{https://supreme.justia.com/cases/federal/us/418/717/}{418 U.S. at 740--41}).

\textbf{Holding vs.\ Dicta:} The holding establishes that absent proof of an inter-district violation producing a significant inter-district segregative effect (or racially drawn district lines), a federal court exceeds its equitable authority by sweeping 53 unoffending suburban school districts into a metropolitan desegregation decree, even if the State itself is derivatively or directly responsible for intra-district violations within the city (\href{https://supreme.justia.com/cases/federal/us/418/717/}{418 U.S. at 744--45, 748--50}). Chief Justice Burger's broader reflections on the virtues of local educational autonomy (citing \textit{San Antonio Independent School District v. Rodriguez}, 411 U.S. 1 (1973)) and his hypothetical catalog of administrative, tax, and bonding problems that consolidation would entail serve as supporting rationale and cautionary dicta rather than a rigid constitutional bar where an inter-district violation \textit{is} proved (\href{https://supreme.justia.com/cases/federal/us/418/717/}{418 U.S. at 741--44}). Likewise, the Court expressly declined to reach the suburban districts' procedural due process claims or state housing-segregation claims that the Sixth Circuit had not relied upon (\href{https://supreme.justia.com/cases/federal/us/418/717/}{418 U.S. at 728 n.7, 752}).

\section{Analysis of Opinions}

\subsection{Majority Opinion (Chief Justice Burger)}
Chief Justice Burger grounded his analysis in two interlocking limitations on federal remedial power. First, he concluded that the lower courts had misapplied \textit{Brown v. Board of Education} and \textit{Swann} by treating racial balance across the metropolitan area as a constitutional mandate (\href{https://supreme.justia.com/cases/federal/us/418/717/}{418 U.S. at 737--41}). Because Judge Roth rejected Detroit-only plans solely on the ground that Detroit's student population was predominantly Black and instead aimed to make no school, grade, or classroom disproportionate to the tri-county racial ratio, the District Court had improperly equated racial imbalance with a constitutional violation (\href{https://supreme.justia.com/cases/federal/us/418/717/}{418 U.S. at 739--41 \& n.19}).

Second, Burger rejected Judge Roth's premise that school district boundaries are mere ``political convenience.'' Declaring that ``[n]o single tradition in public education is more deeply rooted than local control over the operation of schools,'' Burger detailed the extensive statutory authority vested in Michigan's popularly elected local school boards over taxes, bonds, curricula, hiring, and facilities (\href{https://supreme.justia.com/cases/federal/us/418/717/}{418 U.S. at 741--43 \& n.20}). Consolidating 54 independent districts serving over 779,000 pupils into a super-district would force the federal district judge to act as a \textit{de facto} ``legislative authority'' and area-wide ``school superintendent''---roles judges are ill-equipped to perform (\href{https://supreme.justia.com/cases/federal/us/418/717/}{418 U.S. at 743--44}).

Applying \textit{Swann}'s rule that the violation bounds the remedy, Burger held that boundary lines may be crossed only upon proof of an inter-district violation and inter-district segregative effect (\href{https://supreme.justia.com/cases/federal/us/418/717/}{418 U.S. at 744--45}). Even assuming \textit{arguendo} that Michigan was derivatively responsible for Detroit's violations, the record showed no cross-district segregative impact from Act 48, state construction oversight, or funding disparities, and a single late-1950s tuition arrangement between the tiny Carver School District and Detroit could not justify sweeping 52 uninvolved suburban districts into a metropolitan decree (\href{https://supreme.justia.com/cases/federal/us/418/717/}{418 U.S. at 748--52}).

\subsection{Concurring Opinion (Justice Stewart)}
Justice Stewart wrote separately to emphasize the narrowness of the Court's holding as an exercise of equitable restraint rather than a retreat from \textit{Brown} or \textit{Swann} (\href{https://supreme.justia.com/cases/federal/us/418/717/}{418 U.S. at 753, 757}). He stressed that the case presented no evidence that suburban officials had acted unconstitutionally or that governmental policies caused the demographic divide between a predominantly Black Detroit and predominantly white suburbs---attributing urban-suburban racial disparities instead to ``unknown and perhaps unknowable factors such as in-migration, birth rates, economic changes, or cumulative acts of private racial fears'' (\href{https://supreme.justia.com/cases/federal/us/418/717/}{418 U.S. at 756 n.2}). Crucially, however, Stewart outlined circumstances where inter-district relief \textit{would} be proper: where state officials drew or redrew district lines to separate races, transferred school units between districts, or engaged in ``purposeful, racially discriminatory use of state housing or zoning laws'' (\href{https://supreme.justia.com/cases/federal/us/418/717/}{418 U.S. at 755}).

\subsection{Dissenting Opinion (Justice Douglas)}
Justice Douglas argued that Michigan's public school system was constitutionally and statutorily a unitary state enterprise, not a collection of sovereign local fiefdoms: the State controlled district boundaries, site selection, and school bonds, and local districts were mere state agencies (\href{https://supreme.justia.com/cases/federal/us/418/717/}{418 U.S. at 758--59}). Linking \textit{Milliken} to \textit{San Antonio Independent School District v. Rodriguez}, 411 U.S. 1 (1973), Douglas warned that combining property-tax school finance with a ban on cross-district desegregation created a regime worse than \textit{Plessy v. Ferguson}, leaving Black inner-city students in schools that are ``not only `separate' but `inferior'\,'' (\href{https://supreme.justia.com/cases/federal/us/418/717/}{418 U.S. at 759--61}). Reiterating his concurrence in \href{https://supreme.justia.com/cases/federal/us/413/189/}{\textit{Keyes v. School District No. 1}}, 413 U.S. 189 (1973), Douglas maintained there is no constitutional difference between \textit{de facto} and \textit{de jure} school segregation when state action in housing, zoning, and school siting creates racial ghettos (\href{https://supreme.justia.com/cases/federal/us/418/717/}{418 U.S. at 761--62}).

\subsection{Dissenting Opinion (Justice White, joined by Justices Douglas, Brennan, and Marshall)}
Justice White attacked the majority's rule as an arbitrary barrier erected out of ``undue administrative inconvenience to the State'' (\href{https://supreme.justia.com/cases/federal/us/418/717/}{418 U.S. at 763}). Because the Fourteenth Amendment binds the \textit{State} of Michigan---which possessed plenary authority over local school districts and had directly contributed to Detroit's segregation through Act 48, construction oversight, and transportation-aid discrimination---the State could not insulate itself from effective relief by hiding behind its own administrative subdivisions (\href{https://supreme.justia.com/cases/federal/us/418/717/}{418 U.S. at 763, 770--72, 777--78}). White pointed out that the majority left unchallenged Judge Roth's factual findings that a Detroit-only plan would require 900 new buses, leave many schools 75\% to 90\% Black, and accelerate white flight, whereas a metropolitan plan would require at most 350 new buses, cap rides at 40 minutes, and be ``physically easier and more practicable'' (\href{https://supreme.justia.com/cases/federal/us/418/717/}{418 U.S. at 765--68, 781}). Citing \textit{Brown II}, \textit{Green}, \textit{Swann}, and the reapportionment cases (\textit{Reynolds v. Sims} and \textit{Gomillion v. Lightfoot}), White concluded that federal equity power does not stop at municipal or district lines when crossing them is necessary to eradicate state-imposed segregation (\href{https://supreme.justia.com/cases/federal/us/418/717/}{418 U.S. at 772--79}).

\subsection{Dissenting Opinion (Justice Marshall, joined by Justices Douglas, Brennan, and White)}
Justice Marshall delivered a forceful dissent declaring that, after twenty years of incremental progress since \textit{Brown}, ``the Court today takes a giant step backwards'' (\href{https://supreme.justia.com/cases/federal/us/418/717/}{418 U.S. at 782}; \href{https://en.wikisource.org/wiki/Milliken_v._Bradley/Dissent_Marshall}{Wikisource}). First, he refuted the majority's claim that Judge Roth sought an abstract metropolitan racial balance, demonstrating from the record that Judge Roth's sole objective was curing Detroit's proven \textit{de jure} segregation after concluding that all three Detroit-only plans would leave Detroit's schools racially identifiable and trigger white flight (\href{https://supreme.justia.com/cases/federal/us/418/717/}{418 U.S. at 784--90, 799--803}).

Second, Marshall dismantled the majority's portrayal of Michigan's ``local control,'' marshaling Michigan Supreme Court precedent and statutes showing that Michigan operated a centralized statewide educational system that had already consolidated over 6,700 districts since 1912 without local consent (\href{https://supreme.justia.com/cases/federal/us/418/717/}{418 U.S. at 793--98}). Third, he argued that \textit{de jure} segregation within Detroit inevitably rippled across district boundaries by shaping residential patterns and driving white families to the suburbs, and that confining relief to Detroit allowed the State ``to profit from its own wrong'' (\href{https://supreme.justia.com/cases/federal/us/418/717/}{418 U.S. at 804--06}). Finally, Marshall rejected the majority's administrative-chaos fears---noting that no final plan had even been adopted and that Michigan already used inter-district tuition contracts for vocational and special education---and warned that the decision reflected ``a perceived public mood that we have gone far enough'' rather than neutral constitutional principle (\href{https://supreme.justia.com/cases/federal/us/418/717/}{418 U.S. at 809--14}).

\section{Examples: Future Applications}
\begin{enumerate}
\item \textbf{Same-Side (Inter-District Remedy Barred):} \textit{A federal court finds that an urban school board engaged in intentional attendance-zone gerrymandering within city limits, while five adjacent suburban school districts have neutrally drawn nineteenth-century borders, predominantly white enrollments, and no record of discriminatory acts. Plaintiffs seek cross-district busing because the city district is 80\% Black.} Under \textit{Milliken}, an inter-district remedy is impermissible: without proof of an inter-district violation producing a significant segregative effect across boundaries, federal courts must confine desegregation remedies---including compensatory educational programs later approved in \href{https://supreme.justia.com/cases/federal/us/433/267/}{\textit{Milliken v. Bradley} (\textit{Milliken II}), 433 U.S. 267 (1977)}---to the offending urban district (\href{https://supreme.justia.com/cases/federal/us/418/717/}{418 U.S. at 744--45}).

\item \textbf{Same-Side (Inter-District Remedy Barred):} \textit{To attract white suburban students voluntarily into a formerly \textit{de jure} segregated urban district, a district court orders the State to fund lavish, indefinite salary increases and magnet-school upgrades explicitly designed to improve ``desegregative attractiveness'' across the metropolitan area.} Even without mandatory cross-district busing, a remedy aimed at altering inter-district racial demographics to cure a purely intra-district violation exceeds \textit{Milliken}'s remedial ceiling (a principle later reinforced in \textit{Missouri v. Jenkins}, 515 U.S. 70 (1995); see also \href{https://supreme.justia.com/cases/federal/us/495/33/}{\textit{Missouri v. Jenkins}, 495 U.S. 33 (1990)}).

\item \textbf{Opposite-Side (Inter-District Remedy Permitted):} \textit{After a county-wide school system is ordered to dismantle its statutory dual system, a predominantly white municipality within the county obtains state legislative approval to secede and carve out a new, separate school district whose effect is to preserve white schools and impede desegregation of the remaining county.} Under \textit{Milliken}'s own exceptions (citing \textit{Wright v. Council of the City of Emporia} and \textit{United States v. Scotland Neck Board of Education}), federal courts may enjoin the boundary change or impose cross-district relief because district lines were intentionally manipulated to foster or perpetuate racial segregation (\href{https://supreme.justia.com/cases/federal/us/418/717/}{418 U.S. at 744--45, 755}).

\item \textbf{Opposite-Side (Inter-District Remedy Permitted):} \textit{State education officials and three neighboring suburban districts deliberately maintain a cross-border transfer policy and site all federally subsidized minority public housing exclusively inside the central city while excluding Black transfer students from adjacent suburban schools.} Because racially discriminatory acts by the State and suburban districts directly caused inter-district segregation, plaintiffs satisfy \textit{Milliken}'s threshold of an inter-district violation with significant inter-district effects, authorizing a multi-district remedy (\href{https://supreme.justia.com/cases/federal/us/418/717/}{418 U.S. at 744--45, 755}).

\item \textbf{Fence-Sitter (Close Case):} \textit{A state housing authority and municipal zoning boards in suburban counties enforced racially restrictive zoning and redlining policies during the 1950s and 1960s, while suburban school boards themselves operated facially unitary schools; plaintiffs now sue to compel a metropolitan school desegregation decree.} This is a classic fence-sitter: Chief Justice Burger's majority opinion left housing-segregation claims unaddressed because the Sixth Circuit had not relied on them (\href{https://supreme.justia.com/cases/federal/us/418/717/}{418 U.S. at 728 n.7}), whereas Justice Stewart's pivotal concurrence expressly suggested that ``purposeful, racially discriminatory use of state housing or zoning laws'' could justify an inter-district school decree (\href{https://supreme.justia.com/cases/federal/us/418/717/}{418 U.S. at 755}). Whether plaintiffs prevail turns on proving direct, substantial causation linking intentional state housing discrimination to current inter-district school demographics.
\end{enumerate}

\section{Critique}
\textbf{Originalist, Federalist, and Institutional-Restraint Perspectives:} Defenders of \textit{Milliken} from conservative, federalist, and traditional-equity viewpoints praise the decision for restoring principled boundaries to federal judicial power. Grounded in the maxim that equity cures only proven constitutional wrongs, they argue that imposing sweeping structural decrees on 53 suburban school districts---which were neither found liable nor permitted to litigate liability---violates basic norms of federalism, separation of powers, and democratic accountability (\href{https://supreme.justia.com/cases/federal/us/418/717/}{418 U.S. at 743--45}; \href{https://repository.law.umich.edu/mlr/vol122/iss1/4/}{Mich.\ L.\ Rev.}). Scholars examining the Court's ``tradition of local control'' note that the majority treated local educational autonomy not merely as policy preference, but as an institutional safeguard linked to voting rights, federalism, and limits on Article III remedial authority (\href{https://repository.law.umich.edu/mlr/vol122/iss1/4/}{Mich.\ L.\ Rev.}). Originalist-leaning jurists later invoked \textit{Milliken} to caution that open-ended judicial management of local school systems departs from the historical scope of federal equity jurisdiction (\href{https://scholarlycommons.law.case.edu/cgi/viewcontent.cgi?article=2101&context=caselrev}{Case W.\ Res.\ L.\ Rev.}).

\textbf{Progressive, Historical, and Structural Perspectives:} Progressive constitutional scholars, historians, and civil rights researchers widely regard \textit{Milliken} as the decisive turning point that halted northern school desegregation and entrenched metropolitan racial isolation (\href{https://www.gse.harvard.edu/ideas/ed-magazine/14/06/brown-60-and-milliken-40}{Harvard Graduate School of Education}; \href{https://civilrightsproject.ucla.edu/wp-content/uploads/2024/04/National-Segregation-041624-CORRECTED-for-1.pdf}{UCLA Civil Rights Project}). Critics identify three major flaws. First, doctrinal critics argue that Burger's invocation of ``local control'' ignored Michigan's actual constitutional structure, in which public education is a plenary state function and school districts are state instrumentalities (\href{https://www.cambridge.org/core/journals/history-of-education-quarterly/article/revisiting-milliken-v-bradley-michigan-v-the-city-of-detroit/01AD852371763098FB095AD0F45CA8A5}{History of Education Quarterly}). Second, historians and legal scholars such as Myron Orfield and Gary Orfield emphasize that \textit{Milliken} created an artificial wall between ``\textit{de jure}'' school segregation and ``\textit{de facto}'' residential segregation, overlooking decades of federal, state, and local housing discrimination (such as FHA redlining and restrictive covenants) that produced white suburbs and Black urban cores (\href{https://harvardlawreview.org/wp-content/uploads/2026/03/139-Harv.-L.-Rev.-1089.pdf}{Harv.\ L.\ Rev.}; \href{https://scholarship.law.umn.edu/cgi/viewcontent.cgi?article=1987&context=faculty_articles}{Minn.\ Scholarship Repository}). Third, paired with \textit{San Antonio Independent School District v. Rodriguez} (1973), \textit{Milliken} locked minority urban students into underfunded, racially isolated districts while providing a legal safe harbor for white flight---leaving Detroit's public schools overwhelmingly segregated in the decades that followed (\href{https://www.gse.harvard.edu/ideas/ed-magazine/14/06/brown-60-and-milliken-40}{Harvard Graduate School of Education}; \href{https://civilrightsproject.ucla.edu/wp-content/uploads/2024/04/National-Segregation-041624-CORRECTED-for-1.pdf}{UCLA Civil Rights Project}).

\textbf{Underlying Competing Values:} \textit{Milliken} crystallizes the clash between four fundamental values in Equal Protection jurisprudence: (1) substantive racial integration and effective eradication of state-created caste systems versus formal, perpetrator-centered causality limiting remedies to proven wrongdoers; (2) centralized state accountability under the Fourteenth Amendment versus local democratic self-governance and parental control over neighborhood schools; (3) broad equitable flexibility to vindicate constitutional rights versus institutional limits on federal courts acting as super-legislatures; and (4) the contested boundary between actionable state discrimination and private residential choice.

\section{Key Quotations}
\begin{quote}
``No single tradition in public education is more deeply rooted than local control over the operation of schools; local autonomy has long been thought essential both to the maintenance of community concern and support for public schools and to quality of the educational process.''
\end{quote}
\noindent --- Chief Justice Warren E. Burger, Opinion of the Court, 418 U.S. at 741--42 (\href{https://supreme.justia.com/cases/federal/us/418/717/}{Justia}; \href{https://www.courtlistener.com/opinion/109102/milliken-v-bradley/}{CourtListener}).

\begin{quote}
``Before the boundaries of separate and autonomous school districts may be set aside by consolidating the separate units for remedial purposes or by imposing a cross-district remedy, it must first be shown that there has been a constitutional violation within one district that produces a significant segregative effect in another district. Specifically, it must be shown that racially discriminatory acts of the state or local school districts, or of a single school district have been a substantial cause of inter-district segregation.''
\end{quote}
\noindent --- Chief Justice Warren E. Burger, Opinion of the Court, 418 U.S. at 744--45 (\href{https://supreme.justia.com/cases/federal/us/418/717/}{Justia}; \href{https://www.courtlistener.com/opinion/109102/milliken-v-bradley/}{CourtListener}).

\begin{quote}
``Since the mere fact of different racial compositions in contiguous districts does not itself imply or constitute a violation of the Equal Protection Clause in the absence of a showing that such disparity was imposed, fostered, or encouraged by the State or its political subdivisions, it follows that no inter-district violation was shown in this case.''
\end{quote}
\noindent --- Justice Potter Stewart, Concurring, 418 U.S. at 756 (\href{https://supreme.justia.com/cases/federal/us/418/717/}{Justia}; \href{https://www.courtlistener.com/opinion/109102/milliken-v-bradley/}{CourtListener}).

\begin{quote}
``The core of my disagreement is that deliberate acts of segregation and their consequences will go unremedied not because a remedy would be infeasible or unreasonable in terms of the usual criteria governing school desegregation cases, but because an effective remedy would cause what the Court considers to be undue administrative inconvenience to the State. The result is that the State of Michigan, the entity at which the Fourteenth Amendment is directed, has successfully insulated itself from its duty to provide effective desegregation remedies by vesting sufficient power over its public schools in its local school districts.''
\end{quote}
\noindent --- Justice Byron R. White, Dissenting, 418 U.S. at 763 (\href{https://supreme.justia.com/cases/federal/us/418/717/}{Justia}; \href{https://www.courtlistener.com/opinion/109102/milliken-v-bradley/}{CourtListener}).

\begin{quote}
``We deal here with the right of all of our children, whatever their race, to an equal start in life and to an equal opportunity to reach their full potential as citizens. Those children who have been denied that right in the past deserve better than to see fences thrown up to deny them that right in the future. Our Nation, I fear, will be ill-served by the Court's refusal to remedy separate and unequal education, for unless our children begin to learn together, there is little hope that our people will ever learn to live together.''
\end{quote}
\noindent --- Justice Thurgood Marshall, Dissenting, 418 U.S. at 783 (\href{https://supreme.justia.com/cases/federal/us/418/717/}{Justia}; \href{https://www.courtlistener.com/opinion/109102/milliken-v-bradley/}{CourtListener}).

\end{document}
